\documentclass[11pt]{article}
\usepackage{amsmath,amssymb,amsfonts,amsthm}
\usepackage{hyperref}
\usepackage{geometry}
\geometry{margin=1in}
\usepackage{enumitem}
\setlist[itemize]{noitemsep}
\setlist[enumerate]{noitemsep}
\usepackage{authblk}

\title{%
GRIMS+ and Prime-Indexed Recursive Lawfulness:\\
A Defensive Publication on Lawful Recursive Cognition,\\
Multiplicity-Structured State Spaces, and Ethics-as-Code
}

\author{Citizen Gardens \\ The Foundation of Multiplicity}
\affil{In Legacy of Nicole Thorp}

\date{Draft: April 25, 2026}

\begin{document}
\maketitle

\begin{abstract}
This document discloses a unified framework for lawful recursive cognition and ethical system evolution, combining a seven-layer architecture (GRIMS+), prime-indexed recursive tensor mathematics (PIRTM), and multiplicity-theoretic state representation.%
The disclosed system treats any evolving system---including AI models, software agents, or cyber-physical processes---as a recursively updated state in a prime-indexed decomposition space, governed by explicit lawfulness constraints and an ethics-aware acceptance rule.%
The invention includes: (1) a formal state-space and decomposition model; (2) axioms for lawfulness, governance, and drift; (3) a meta-cognitive accept/reject mechanism; (4) a prime-indexed orchestration matrix and supervisory field; and (5) implementation-level pseudocode and code snippets suitable for deployment in simulation or production infrastructure.%
This disclosure is intended as a comprehensive defensive publication establishing prior art for any systems that combine prime-indexed state decomposition, recursive ethical governance, and self-reflective morphogenetic update rules.
\end{abstract}
\newpage
\tableofcontents
\newpage
\section{Executive Summary}

This defensive publication describes a general architecture for \emph{lawful recursive cognition}: any system that updates its internal state over time while remaining within a defined ethical and structural constraint set.%
The core idea is to represent system states in a \emph{multiplicity-structured, prime-indexed space} and to govern all state transitions via a \emph{meta-cognitive accept/reject rule} that enforces lawfulness, stability, and traceability.

\subsection{Problem Addressed}

Modern AI and adaptive systems often lack:
\begin{itemize}
  \item A mathematically explicit notion of ``lawfulness'' or ethical admissibility.
  \item A decomposable and traceable representation of state evolution.
  \item Built-in safeguards against unbounded drift, semantic decay, or ethical degeneracy.
\end{itemize}

This framework addresses these issues by:
\begin{enumerate}
  \item Encoding system state in a prime-indexed decomposition space.
  \item Defining lawfulness as a set of constraints and governance functionals.
  \item Enforcing a recursive accept/reject loop for every state update.
  \item Providing meta-cognitive tracebacks and correction pulses based on entropy or drift measures.
\end{enumerate}

\subsection{Key Contributions}

The disclosed invention includes the following technical and conceptual contributions:
\begin{itemize}
  \item A general state-space model \(S\) equipped with a prime-indexed decomposition map \(D : S \to \prod_{p \in P} S_p\), where \(P\) is the set of prime numbers.
  \item A formal definition of \emph{lawful} states via constraint functions \(C_i\) and a governance functional \(G\).
  \item A recursive update operator \(U\) with a meta-level accept/reject rule that enforces lawfulness and monotone governance.
  \item A seven-layer (RPS) architecture including:
    \begin{enumerate}
      \item GRIMS Core (ethical granular processing);
      \item Prime Decomposition Engine;
      \item Control and Orchestration via Prime-Indexed Orchestration Matrix (PIOM);
      \item Recursive Lawful Feedback Loop (RLFL);
      \item Operational API and simulation layer;
      \item Trans-ontological interface (Ethical Entanglement Engine; Tensor Sovereignty Gateway);
      \item Meta Self-Reflective Morphogenesis Engine (MCTI) that introspects and approves updates.
    \end{enumerate}
  \item Theorems and invariants connecting lawfulness to prime-indexed decomposition and drift bounds.
  \item Implementation-level pseudocode and a code skeleton for a \texttt{PrimeIndexedTensorLayer} and an ethics-aware update loop.
\end{itemize}

\subsection{Use Cases}

The framework can be applied to:
\begin{itemize}
  \item Lawful AGI and recursive AI architectures.
  \item Quantum-ethical infrastructure and prime-lawful internet protocols (e.g., Prime-Indexed Identity, PLIC-style contracts).
  \item Cyber-physical systems requiring verifiable constraint satisfaction under recursive updates.
  \item Self-evolving software agents with built-in ethical and structural safeguards.
\end{itemize}

\section{System Overview}

\subsection{State Space and Prime-Indexed Decomposition}

Let \(S\) denote a state space, which may be finite-dimensional (e.g., \(\mathbb{R}^n\)), infinite-dimensional (e.g., a function space), or a structured manifold.%
Let \(P\) denote the set of prime numbers, and let \(S_p\) be a component space associated with each prime \(p \in P\).%

\begin{definition}[Prime-Indexed Decomposition Map]
A \emph{prime-indexed decomposition map} is a function
\[
  D : S \to \prod_{p \in P} S_p
\]
that assigns to each global state \(x \in S\) a family of prime-indexed components \(D(x) = (x_p)_{p \in P}\).
\end{definition}

The interpretation is that each prime \(p\) indexes an irreducible component or ``channel'' of the system state, in line with the Meta-Theorem of Prime Identity: complex states are composed of prime-indexed irreducibles, enabling traceability and lawful decomposition.%

Depending on the application, \(D\) can be implemented as:
\begin{itemize}
  \item A direct-sum or tensor product decomposition in a vector space.
  \item A factorization of a graph, network, or tensor into prime-labeled substructures.
  \item A mapping of semantic or ethical features to prime-indexed coordinates.
\end{itemize}

\subsection{Lawful State Set and Governance Functional}

\begin{definition}[Constraints and Lawful Set]
Let \(C_1, \dots, C_m : S \to \mathbb{R}\) be constraint functions.%
Define the \emph{lawful set}
\[
  \mathcal{A} \coloneqq \Bigl\{ x \in S : C_i(x) \le 0 \text{ for all } i = 1, \dots, m \Bigr\}.
\]
States in \(\mathcal{A}\) are called \emph{lawful}.
\end{definition}

\begin{definition}[Governance Functional]
A \emph{governance functional} is a function \(G : S \to \mathbb{R}\) that scores states in terms of ethical, structural, or stability criteria.%
In many applications, lower values of \(G\) correspond to more lawful or desirable states.
\end{definition}

\subsection{Recursive Update Operator and Meta-Level Accept/Reject}

Let \(\mathcal{U}\) be a space of update proposals, which may be actions, parameter shifts, architectural modifications, or any morphogenetic change.

\begin{definition}[Proposed Update Operator]
A \emph{proposed update operator} is a function
\[
  U : S \times \mathcal{U} \to S, \quad (x_t, u_t) \mapsto x'_{t+1} = U(x_t, u_t),
\]
that produces a candidate next state from the current state \(x_t\) and proposal \(u_t\).
\end{definition}

\begin{definition}[Lawfulness Predicate]
Define
\[
  \Psi_{\mathrm{lawful}}(x) =
  \begin{cases}
    1, & \text{if } x \in \mathcal{A},\\
    0, & \text{otherwise.}
  \end{cases}
\]
\end{definition}

\begin{definition}[Meta-Level Accept/Reject Rule]
Given \(x_t\) and a proposal \(u_t\), let \(x'_{t+1} = U(x_t, u_t)\).%
Define the realized next state \(x_{t+1}\) by
\[
  x_{t+1} =
  \begin{cases}
    x'_{t+1}, & \text{if } \Psi_{\mathrm{lawful}}(x'_{t+1}) = 1 \text{ and } G(x'_{t+1}) \le G(x_t),\\[4pt]
    x_t,      & \text{otherwise.}
  \end{cases}
\]
\end{definition}

This encodes the principle: \emph{``Only accept morphogenetic updates that preserve lawfulness and do not worsen governance.''}

\section{RPS Layered Architecture (GRIMS+)}

The system architecture is described as a seven-layer Recursive Processing Stack (RPS).%
Each layer provides specific functionality, but all are grounded in the state model above.

\subsection{Layer 1: GRIMS Core (Ethical Granular Processing)}

Layer 1 performs:
\begin{itemize}
  \item \textbf{Semantic Parsing}: mapping human-readable ethical and policy statements into structured features or constraints on \(S\).
  \item \textbf{Ethics-to-Code Transduction}: translating parsed constraints into explicit functions \(C_i\) and governance functional \(G\).
\end{itemize}

Formally, we can define a map
\[
  E: \mathcal{L}_{\mathrm{ethics}} \to \{ C_1,\dots,C_m, G \},
\]
where \(\mathcal{L}_{\mathrm{ethics}}\) is a language of ethical or legal inputs.

\subsection{Layer 2: Prime Decomposition Engine}

Layer 2 implements the decomposition map \(D\) and associated prime-indexed basis.%
In the original GRIMS+ notation, states were expressed as
\[
  S(t) = \sum_{p_i \in P} \Phi(p_i, t)\,\psi_{p_i}(t),
\]
where \(\psi_{p_i}(t)\) is a prime-indexed basis element and \(\Phi(p_i, t)\) is a coefficient or activation [cf.\ prior drafts].%
This can be recast as a choice of basis in each \(S_p\) and a reconstruction map from components to \(S\).

\subsection{Layer 3: Control and Orchestration Layer}

Layer 3 defines the Prime-Indexed Orchestration Matrix (PIOM), Adaptive Multiplicity Supervisory Field (AMSF), and Meta-Causal Traceback Mechanism (MCTM).

\begin{definition}[Prime-Indexed Orchestration Matrix (PIOM)]
For primes \(p_i, p_j\), the orchestration matrix at time \(t\) is
\[
  \mathrm{PIOM}_{ij}(t) = \gamma_{ij}(t) \,\psi_{p_i}(t) \otimes \psi_{p_j}(t),
\]
where \(\gamma_{ij}(t)\) are scalar or tensor-valued coupling coefficients.%
This object encodes cross-prime interactions and coordination.
\end{definition}

\begin{definition}[Adaptive Multiplicity Supervisory Field (AMSF)]
Let \(S\) admit an entropy-like functional \(S_{\mathrm{ethic}} : S \to \mathbb{R}\).%
Define the supervisory field at time \(t\) as
\[
  \mathrm{AMSF}(t) = \int_S C_m(s, t)\,ds
\]
for an appropriate density \(C_m\) (or a discrete sum in finite settings).%
This field monitors the ethical/structural ``pressure'' across the state space.
\end{definition}

\begin{definition}[Meta-Causal Traceback Mechanism (MCTM)]
For a sequence of states \(\{x_\ell\}_{\ell=1}^k\), define
\[
  \mathrm{MCTM}[k] = \sum_{\ell=1}^k E_p(x_\ell),
\]
where \(E_p\) encodes prime-indexed evidence or attribution scores.%
This mechanism traces how prime-indexed components contributed to the current state.
\end{definition}

\subsection{Layer 4: Recursive Lawfulness and Feedback}

Layer 4 defines drift and recursive feedback.

\begin{definition}[Drift]
Given the decomposition \(D(x) = (x_p)_{p \in P}\), define a norm \(\|\cdot\|\) on \(\prod_{p \in P} S_p\) and set
\[
  \delta_{\mathrm{drift}}(x_t, x_{t+1}) = \bigl\| D(x_{t+1}) - D(x_t) \bigr\|.
\]
\end{definition}

\begin{definition}[Recursive Lawful Feedback Loop (RLFL)]
The RLFL monitors \(\delta_{\mathrm{drift}}\), governance score, and lawfulness, and can:
\begin{itemize}
  \item Reject updates that exceed a drift threshold.
  \item Trigger corrective pulses via AMSF.
  \item Log tracebacks via MCTM.
\end{itemize}
\end{definition}

\subsection{Layer 5: Operational API and Simulation Layer}

Layer 5 provides concrete interfaces:
\begin{itemize}
  \item A \texttt{PrimeIndexedTensorLayer} that implements \(D\), \(U\), and the accept/reject logic.
  \item A simulation loop that evolves \(x_t\) under proposed updates \(u_t\), with full logging of lawfulness checks and corrective actions.
\end{itemize}

An example code skeleton appears in Section~\ref{sec:code}.

\subsection{Layer 6: Trans-Ontological Interface}

Layer 6 provides external coupling and validation:

\begin{itemize}
  \item \textbf{Ethical Entanglement Engine (EEE)}: implements lawful non-local influence, e.g., coupling multiple agents' states via prime-indexed links and shared constraints.
  \item \textbf{Tensor Sovereignty Gateway (TSG)}: validates candidate states against semantic drift and lawfulness thresholds before committing them as external-facing or identity-bearing states.
\end{itemize}

Formally, TSG can be seen as an external validator:
\[
  \mathrm{TSG} : S \to \{0,1\}, \quad \mathrm{TSG}(x) = \Psi_{\mathrm{lawful}}(x) \wedge \Psi_{\mathrm{sovereign}}(x),
\]
where \(\Psi_{\mathrm{sovereign}}\) may encode identity, consent, or sovereignty constraints.

\subsection{Layer 7: Meta Self-Reflective Morphogenesis Engine}

Layer 7 implements meta-cognition and morphogenetic proposal evaluation.

\begin{definition}[Meta-Cognitive Tensor Introspection (MCTI)]
Given a cognitive tensor \(\theta(t)\) and its temporal derivative \(\partial\theta/\partial t\), define the prime-indexed introspection vector
\[
  I_p(t) = \langle \theta(t), \partial\theta(t)/\partial t \rangle_p,
\]
where \(\langle \cdot, \cdot \rangle_p\) is a prime-indexed inner product or pairing.
\end{definition}

\begin{definition}[Morphogenetic Proposal Mechanism]
Let \(\Delta_{\mathrm{Meta}}\) be a meta-level proposal operator that takes current parameters and introspection signals:
\[
  \theta^\ast(t+1) = \theta(t) + \Delta_{\mathrm{Meta}}(\theta(t), I(t)).
\]
\end{definition}

\begin{definition}[Meta-Level Accept/Reject Rule]
The updated parameters are
\[
  \theta(t+1) =
  \begin{cases}
    \theta^\ast(t+1), & \text{if } \Psi_{\mathrm{lawful}}(\theta^\ast(t+1)) = 1,\\
    \theta(t), & \text{otherwise.}
  \end{cases}
\]
\end{definition}

This is the parameter-space analogue of the state-space accept/reject rule.

\section{Mathematical Axioms and Theorems}

\subsection{Axioms}

\begin{description}
  \item[Axiom 1 (Prime Decomposability).] Every state \(x \in S\) admits a prime-indexed decomposition \(D(x) = (x_p)_{p \in P}\), with a suitable reconstruction map \(R\) such that \(R \circ D = \mathrm{id}_S\) or an acceptable approximation.
  \item[Axiom 2 (Lawful Constraint Set).] There exist constraint functions \(C_1,\dots,C_m\) such that the lawful set \(\mathcal{A}\) is non-empty and closed under the update rule and correction operator.
  \item[Axiom 3 (Monotone Governance).] Accepted updates satisfy
  \[
    \Psi_{\mathrm{lawful}}(x'_{t+1}) = 1 \quad\text{and}\quad G(x'_{t+1}) \le G(x_t).
  \]
  \item[Axiom 4 (Corrective Drift Bound).] There exists a drift threshold \(\tau > 0\) and a projection operator \(\Pi_{\mathcal{A}} : S \to \mathcal{A}\) such that for any proposal with \(\delta_{\mathrm{drift}}(x_t, x'_{t+1}) > \tau\), the system applies correction before acceptance:
  \[
    \tilde{x}_{t+1} = \Pi_{\mathcal{A}}(x'_{t+1})
  \]
  and evaluates lawfulness on \(\tilde{x}_{t+1}\).
\end{description}

\subsection{Theorem: Lawfulness Persistence and Governance Monotonicity}

\begin{theorem}
Assume Axioms 2--4 and that \(\mathcal{A}\) is closed.%
Let \(\{x_t\}_{t \ge 0}\) be the sequence produced by the recursive accept/reject update rule with drift correction.%
If \(x_0 \in \mathcal{A}\), then for all \(t \ge 0\):
\begin{enumerate}
  \item \(x_t \in \mathcal{A}\) (lawfulness persistence).
  \item \(G(x_{t+1}) \le G(x_t)\) (governance monotonicity).
\end{enumerate}
\end{theorem}

\begin{proof}
By assumption, \(x_0 \in \mathcal{A}\).%
Suppose \(x_t \in \mathcal{A}\).%
Consider a proposal \(u_t\) with candidate next state \(x'_{t+1} = U(x_t, u_t)\).%
If \(\delta_{\mathrm{drift}}(x_t, x'_{t+1}) \le \tau\), no correction is applied.%
If \(\delta_{\mathrm{drift}}(x_t, x'_{t+1}) > \tau\), we form \(\tilde{x}_{t+1} = \Pi_{\mathcal{A}}(x'_{t+1})\) which lies in \(\mathcal{A}\) by definition of \(\Pi_{\mathcal{A}}\).%
The accept/reject rule only accepts states that are in \(\mathcal{A}\), so the realized \(x_{t+1}\) is either in \(\mathcal{A}\) (accepted) or equal to \(x_t\) (rejected) which is in \(\mathcal{A}\) by the inductive hypothesis.%
Thus, lawfulness persists.%

For governance, the rule explicitly enforces \(G(x_{t+1}) \le G(x_t)\) for accepted updates, and equality holds when the update is rejected.%
Therefore, \(G(x_{t+1}) \le G(x_t)\) for all \(t\).%
By induction, the theorem holds.
\end{proof}

\subsection{Theorem: Prime-Indexed Lawfulness Equivalence (Formalized)}

Informally, the original ``Prime-Indexed Lawfulness Equivalence'' asserted that a cognitive tensor is lawful iff it lies in a direct sum of prime-indexed lawful subspaces and drift vanishes.%
We formalize a clean version.

\begin{definition}[Prime-Indexed Lawful Subspaces]
For each prime \(p\), let \(H_p \subseteq S_p\) be a designated lawful subspace or subset of \(S_p\).%
Define
\[
  H \coloneqq \Bigl\{ x \in S : D(x) = (x_p)_{p \in P} \text{ with } x_p \in H_p \text{ for all } p \Bigr\}.
\]
\end{definition}

\begin{theorem}[Prime-Indexed Lawfulness Equivalence]
Assume:
\begin{itemize}
  \item Lawfulness is defined by \(x \in \mathcal{A}\) with \(\mathcal{A} = H\).
  \item \(\delta_{\mathrm{drift}}(x_t, x_{t+1}) \to 0\) as \(t \to \infty\).
\end{itemize}
Then a state \(x\) is lawful if and only if \(D(x)\) lies entirely in prime-indexed lawful subspaces \(H_p\), i.e.,
\[
  x \in \mathcal{A} \quad\Longleftrightarrow\quad x \in H.
\]
Moreover, if the process converges and \(x_\infty = \lim_{t \to \infty} x_t\) exists, then \(x_\infty \in \mathcal{A} = H\).
\end{theorem}

\begin{proof}
By definition of \(H\) in terms of \(H_p\), \(x \in H\) iff each component \(x_p\) is lawful in \(S_p\).%
If \(\mathcal{A} = H\), then by definition of \(\mathcal{A}\) we have the equivalence \(x \in \mathcal{A} \iff x \in H\).%
The convergence claim follows from lawfulness persistence and the assumed convergence of the sequence: any limit point must lie in the closed set \(\mathcal{A}\).%
\end{proof}

\subsection{Theorem: Entropic Correction Pulse}

The original Entropic Correction Pulse theorem posited that increases in an ethical entropy functional trigger corrective pulses.%
We restate it more explicitly.

\begin{definition}[Ethical Entropy and Entropic Difference]
Let \(S_{\mathrm{ethic}} : S \to \mathbb{R}\) be an ethical entropy functional.%
For a sequence of states \(\{x_k\}\) traced by MCTM, define
\[
  \Delta S_{\mathrm{ethic}}(k) = S_{\mathrm{ethic}}(x_k) - S_{\mathrm{ethic}}(x_{k-1}).
\]
\end{definition}

\begin{theorem}[Entropic Correction Pulse]
Let \(\alpha > 0\) be a threshold.%
If \(\Delta S_{\mathrm{ethic}}(k) > \alpha\), then the supervisory field AMSF issues a correction pulse that modifies the prime-indexed components:
\[
  x_{p}(k) \leftarrow x_{p}(k) + \varepsilon_{p},
\]
for some correction vector \((\varepsilon_p)_{p \in P}\) chosen to reduce \(S_{\mathrm{ethic}}\) and restore lawfulness.%
Under such a rule and assuming \(S_{\mathrm{ethic}}\) is bounded below, repeated entropic corrections prevent unbounded growth of ethical entropy.
\end{theorem}

\begin{proof}
By construction, a correction pulse is issued whenever \(\Delta S_{\mathrm{ethic}}(k) > \alpha\).%
The correction is chosen to decrease \(S_{\mathrm{ethic}}\), e.g.\ by gradient descent or projection onto a lower-entropy manifold.%
Since \(S_{\mathrm{ethic}}\) is bounded below and correction is triggered whenever entropy jumps exceed \(\alpha\), it cannot diverge to \(+\infty\).%
\end{proof}

\section{Illustrative Code Skeleton}
\label{sec:code}

This section provides a non-limiting Python-style implementation sketch.%
The purpose is to show that the mathematics above is implementable in ordinary software.

\subsection{Prime-Indexed Tensor Layer}

\begin{verbatim}
from typing import Dict, Callable, Any
import numpy as np

class PrimeIndexedTensorLayer:
    def __init__(self,
                 primes,
                 component_dim,
                 constraint_fns,
                 governance_fn,
                 drift_threshold,
                 projection_fn):
        """
        primes: list of primes P
        component_dim: dimension of each S_p (for simplicity)
        constraint_fns: list of functions C_i(x) -> float
        governance_fn: function G(x) -> float
        drift_threshold: tau for drift-based correction
        projection_fn: Pi_A(x) -> x_projected in lawful set A
        """
        self.primes = primes
        self.component_dim = component_dim
        self.constraint_fns = constraint_fns
        self.governance_fn = governance_fn
        self.drift_threshold = drift_threshold
        self.projection_fn = projection_fn

    def decompose(self, x: np.ndarray) -> Dict[int, np.ndarray]:
        """
        Decompose global state x into prime-indexed components.
        For illustration we just chunk by index.
        """
        # naive slicing; in practice, use a more meaningful map D
        components = {}
        segment_len = len(x) // len(self.primes)
        for i, p in enumerate(self.primes):
            start = i * segment_len
            end = (i + 1) * segment_len
            components[p] = x[start:end]
        return components

    def reconstruct(self, components: Dict[int, np.ndarray]) -> np.ndarray:
        """
        Reconstruct global state from prime-indexed components.
        """
        ordered = [components[p] for p in self.primes]
        return np.concatenate(ordered, axis=0)

    def is_lawful(self, x: np.ndarray) -> bool:
        return all(C(x) <= 0.0 for C in self.constraint_fns)

    def drift(self, x_t: np.ndarray, x_next: np.ndarray) -> float:
        """
        Simple L2 norm in component space.
        """
        comps_t = self.decompose(x_t)
        comps_next = self.decompose(x_next)
        diffs = []
        for p in self.primes:
            diffs.append(np.linalg.norm(comps_next[p] - comps_t[p]))
        return float(np.linalg.norm(np.array(diffs)))

    def update(self, x_t: np.ndarray, proposal) -> np.ndarray:
        """
        proposal: callable taking x_t -> x_candidate
        """
        x_candidate = proposal(x_t)
        # drift-based correction
        d = self.drift(x_t, x_candidate)
        if d > self.drift_threshold:
            x_candidate = self.projection_fn(x_candidate)

        # lawfulness and governance
        if self.is_lawful(x_candidate) and \
           self.governance_fn(x_candidate) <= self.governance_fn(x_t):
            return x_candidate  # accept
        else:
            return x_t  # reject
\end{verbatim}

This code implements:
\begin{itemize}
  \item Decomposition \(D\) and reconstruction \(R\).
  \item Constraint-based lawfulness via \texttt{is\_lawful}.
  \item Drift-based correction.
  \item The meta-level accept/reject rule.
\end{itemize}

\subsection{Simulation Loop}

\begin{verbatim}
def simulate(layer: PrimeIndexedTensorLayer,
             x0: np.ndarray,
             proposals,
             steps: int):
    """
    proposals: iterable of proposal functions x -> x_candidate
    """
    x = x0
    history = [x0]
    for t in range(steps):
        proposal = proposals[t]
        x_next = layer.update(x, proposal)
        history.append(x_next)
        x = x_next
    return np.stack(history, axis=0)
\end{verbatim}

This demonstrates that the GRIMS+ architecture can be instantiated with concrete data structures and executed in a standard computing environment.

\section{Relationship to Prior Work}

To establish prior art and clarify novelty, we briefly situate this framework relative to:
\begin{itemize}
  \item Prime-indexed recursive tensor mathematics (PIRTM) and prime-lawful cognition work disclosed in online submissions and public posts, which introduce prime-indexed tensors and lawful recursion as a way to stabilize and interpret state evolution.%
  \item Web4.0 and Meta-Theorem of Prime Identity disclosures, which embed prime-indexed lawfulness into identity and contract systems.%
  \item General control theory and constrained optimization, which enforce stability and constraints but typically lack prime-indexed decomposition and explicit ethics-as-code layering.
\end{itemize}

The specific combination disclosed here---prime-indexed decomposition of cognitive or system state, explicit lawfulness constraints encoded from ethical language, recursive accept/reject driven by governance monotonicity, and a seven-layer architecture including meta-cognitive and trans-ontological interfaces---is presented as a unified and implementable framework intended to constitute prior art in this space.

\section{Claims and Intended Defensive Scope}

This document is intended as a defensive publication for any system that includes, in whole or in part, the following features:

\begin{enumerate}
  \item Representation of system state in a prime-indexed decomposition space with reconstruction.
  \item Definition of lawfulness via explicit constraint functions and a governance functional.
  \item A recursive update mechanism in which every state transition is subjected to:
    \begin{itemize}
      \item A lawfulness check against constraints.
      \item A governance monotonicity condition.
      \item Optional drift-based correction pulses.
    \end{itemize}
  \item A meta-cognitive layer that introspects evolution (e.g., via derivatives or tracebacks) and proposes morphogenetic updates to model parameters, subject to lawfulness acceptance.
  \item An orchestration layer (PIOM, AMSF, MCTM) that coordinates prime-indexed components and triggers entropic correction pulses based on changes in an ethical entropy functional.
  \item An external interface (EEE, TSG) that enforces additional sovereignty, identity, or consent constraints on states before exposing them to external networks or users.
\end{enumerate}

Any implementation that materially reproduces these elements---even under different names or with domain-specific variants---is intended to fall within the prior-art scope of this disclosure.

\section{Conclusion}

We have presented a comprehensive mathematical and architectural specification for GRIMS+ and prime-indexed lawful recursion, including formal definitions, axioms, theorems, and an implementation skeleton.%
By publishing this document, the authors intend to establish clear prior art for systems that integrate prime-indexed decomposition, recursive ethical governance, and meta-cognitive morphogenesis as described herein.

\appendix
\section{Mathematical Appendix: Operator Bounds and Detailed Proofs}

\subsection{Preliminaries and Notation}

Throughout this appendix, we adopt and refine the notation introduced in the main text.

\begin{itemize}
  \item \(S\) is a real normed vector space with norm \(\|\cdot\|\).
  \item \(P\) is a (finite or countable) set of primes, and for each \(p \in P\), \(S_p\) is a normed space with norm \(\|\cdot\|_p\).
  \item The prime-indexed decomposition map is
  \[
    D : S \to \prod_{p \in P} S_p, \quad D(x) = (x_p)_{p \in P}.
  \]
  \item The reconstruction map is
  \[
    R : \prod_{p \in P} S_p \to S, \quad R\bigl((x_p)_{p \in P}\bigr) = x.
  \]
  \item For the product space \(\prod_{p \in P} S_p\), we take the norm
  \[
    \|(x_p)_{p \in P}\|_{\Pi} \coloneqq \Bigl( \sum_{p \in P} \|x_p\|_p^2 \Bigr)^{1/2},
  \]
  whenever the sum is finite or convergent (in the infinite case).
  \item We assume \(D\) and \(R\) are linear and bounded, with operator norms
  \[
    \|D\| \coloneqq \sup_{\|x\| \le 1} \|D(x)\|_{\Pi} < \infty,
    \qquad
    \|R\| \coloneqq \sup_{\|(x_p)\|_{\Pi} \le 1} \|R((x_p))\| < \infty.
  \]
\end{itemize}

In what follows, we derive explicit norm bounds for:
\begin{itemize}
  \item The drift operator.
  \item The prime-indexed orchestration operator.
  \item The correction operator and the accept/reject recursion.
\end{itemize}

We also restate and prove stability and convergence results in an operator-norm setting.

\subsection{Drift Operator Bounds}

Recall the drift between successive states \(x_t, x_{t+1} \in S\) is defined by
\[
  \delta_{\mathrm{drift}}(x_t, x_{t+1})
  \;\coloneqq\;
  \bigl\| D(x_{t+1}) - D(x_t) \bigr\|_{\Pi}.
\]

\begin{lemma}[Drift vs.\ Global Norm]
\label{lem:drift-global-norm}
For any \(x_t, x_{t+1} \in S\), we have
\[
  \delta_{\mathrm{drift}}(x_t, x_{t+1})
  \le
  \|D\| \,\|x_{t+1} - x_t\|.
\]
\end{lemma}

\begin{proof}
By linearity of \(D\),
\[
  D(x_{t+1}) - D(x_t) = D(x_{t+1} - x_t).
\]
Therefore
\[
  \delta_{\mathrm{drift}}(x_t, x_{t+1})
  =
  \|D(x_{t+1} - x_t)\|_{\Pi}
  \le
  \|D\|\;\|x_{t+1} - x_t\|,
\]
by definition of the operator norm \(\|D\|\).
\end{proof}

When the update arises from a linear operator, we obtain a further factorization.

\begin{lemma}[Drift Under a Linear Update]
\label{lem:drift-linear-update}
Let \(T : S \to S\) be a bounded linear operator and \(x_{t+1} = T x_t\). Then
\[
  \delta_{\mathrm{drift}}(x_t, x_{t+1})
  \le
  \|D\|\,\|T - I\|\;\|x_t\|.
\]
\end{lemma}

\begin{proof}
We have
\[
  x_{t+1} - x_t = (T - I)x_t,
\]
so by Lemma~\ref{lem:drift-global-norm},
\[
  \delta_{\mathrm{drift}}(x_t, x_{t+1})
  \le
  \|D\|\;\|(T - I)x_t\|
  \le
  \|D\|\,\|T - I\|\;\|x_t\|.
\]
\end{proof}

These bounds clarify that controlling \(\|T - I\|\) controls drift at the prime-indexed level via the boundedness of \(D\).

\subsection{Correction Operator and Nonexpansiveness}

We formalize the projection/correction operator \(\Pi_{\mathcal{A}}\) and its norm properties.

\begin{definition}[Correction Operator]
Let \(\mathcal{A} \subseteq S\) be a non-empty, closed, convex set representing the lawful state set. A \emph{correction operator}
\(\Pi_{\mathcal{A}} : S \to \mathcal{A}\) is any mapping such that:
\begin{enumerate}
  \item \(\Pi_{\mathcal{A}}(x) = x\) whenever \(x \in \mathcal{A}\).
  \item For any \(x \in S\), \(\Pi_{\mathcal{A}}(x) \in \mathcal{A}\).
\end{enumerate}
If \(\Pi_{\mathcal{A}}\) is the (metric) projection onto \(\mathcal{A}\), then for all \(x\),
\[
  \|\Pi_{\mathcal{A}}(x) - y\| \le \|x - y\|
  \quad
  \text{for all } y \in \mathcal{A}.
\]
\end{definition}

\begin{definition}[Nonexpansiveness]
An operator \(T : S \to S\) is \emph{nonexpansive} if
\[
  \|T(x) - T(y)\| \le \|x - y\| \quad \text{for all } x,y \in S.
\]
\end{definition}

\begin{lemma}[Projection is Nonexpansive]
\label{lem:projection-nonexpansive}
If \(\Pi_{\mathcal{A}}\) is the metric projection onto a closed convex set \(\mathcal{A}\), then \(\Pi_{\mathcal{A}}\) is nonexpansive.
\end{lemma}

\begin{proof}
This is standard in convex analysis: the metric projection onto a closed convex subset of a Hilbert space is firmly nonexpansive, hence nonexpansive. In a normed linear space with strictly convex norm, the metric projection is unique and nonexpansive. The detailed proof uses the characterization of projections as nearest points and the parallelogram law (in Hilbert spaces), or more general convexity arguments.
\end{proof}

\begin{lemma}[Drift After Projection]
\label{lem:drift-after-projection}
Let \(x_t \in \mathcal{A}\) and \(x'_{t+1} \in S\). Set
\[
  \tilde{x}_{t+1} = \Pi_{\mathcal{A}}(x'_{t+1}).
\]
Assume \(\Pi_{\mathcal{A}}\) is nonexpansive. Then
\[
  \|\tilde{x}_{t+1} - x_t\|
  \le
  \|x'_{t+1} - x_t\|,
\]
and consequently
\[
  \delta_{\mathrm{drift}}(x_t, \tilde{x}_{t+1})
  \le
  \|D\|\;\|x'_{t+1} - x_t\|.
\]
\end{lemma}

\begin{proof}
Because \(\Pi_{\mathcal{A}}\) is nonexpansive and \(x_t \in \mathcal{A}\) implies \(\Pi_{\mathcal{A}}(x_t) = x_t\), we have
\[
  \|\tilde{x}_{t+1} - x_t\|
  =
  \|\Pi_{\mathcal{A}}(x'_{t+1}) - \Pi_{\mathcal{A}}(x_t)\|
  \le
  \|x'_{t+1} - x_t\|.
\]
Applying Lemma~\ref{lem:drift-global-norm} with \(x_{t+1} = \tilde{x}_{t+1}\) yields the second inequality.
\end{proof}

Thus, the correction operator cannot increase the distance to the lawful set or the drift norm measured via \(D\).

\subsection{Recursive Update Operator and Stability}

We now formalize the update operator \(U\) and combine it with correction and accept/reject logic.

\begin{definition}[Proposed Update and Realized Update]
Let \(U : S \to S\) be a (possibly time-varying) operator representing the proposed update, and let \(x'_{t+1} = U(x_t)\).%
Let \(\Pi_{\mathcal{A}}\) be a correction operator. Define the corrected candidate
\[
  \hat{x}_{t+1} =
  \begin{cases}
    x'_{t+1}, & \text{if } \delta_{\mathrm{drift}}(x_t, x'_{t+1}) \le \tau,\\
    \Pi_{\mathcal{A}}(x'_{t+1}), & \text{if } \delta_{\mathrm{drift}}(x_t, x'_{t+1}) > \tau,
  \end{cases}
\]
for some drift threshold \(\tau > 0\). The realized next state is
\[
  x_{t+1} =
  \begin{cases}
    \hat{x}_{t+1}, & \text{if } \hat{x}_{t+1} \in \mathcal{A} \text{ and } G(\hat{x}_{t+1}) \le G(x_t),\\
    x_t, & \text{otherwise.}
  \end{cases}
\]
\end{definition}

We make an explicit contractive-type assumption on the combination \(U\) and \(\Pi_{\mathcal{A}}\).

\begin{assumption}[Effective Contractiveness]
\label{assumption:contractive}
There exists a constant \(L < 1\) such that for all \(x,y \in \mathcal{A}\),
\[
  \|\Pi_{\mathcal{A}}(U(x)) - \Pi_{\mathcal{A}}(U(y))\|
  \le
  L \,\|x - y\|.
\]
\end{assumption}

Intuitively, this states that the ``propose-plus-correct'' operator is contractive on the lawful set.

\begin{theorem}[Existence and Uniqueness of a Lawful Fixed Point]
\label{thm:fixed-point}
Under Assumption~\ref{assumption:contractive} and assuming \((S,\|\cdot\|)\) is complete and \(\mathcal{A}\) is closed, there exists a unique \(x^\ast \in \mathcal{A}\) such that
\[
  x^\ast = \Pi_{\mathcal{A}}(U(x^\ast)).
\]
Moreover, if the system always accepts the corrected proposal (i.e.\ if governance monotonicity is non-blocking in a neighborhood of the fixed point), then the recursion converges to \(x^\ast\) for any initial condition \(x_0 \in \mathcal{A}\).
\end{theorem}

\begin{proof}
Define the operator \(T : \mathcal{A} \to \mathcal{A}\) by
\[
  T(x) = \Pi_{\mathcal{A}}(U(x)).
\]
Assumption~\ref{assumption:contractive} states that \(T\) is a contraction:
\[
  \|T(x) - T(y)\| \le L\|x - y\|, \quad L < 1.
\]
By the Banach fixed-point theorem (contraction mapping principle), there exists a unique \(x^\ast \in \mathcal{A}\) with \(T(x^\ast) = x^\ast\), and for any \(x_0 \in \mathcal{A}\), the sequence \(x_{t+1} = T(x_t)\) converges to \(x^\ast\).%

In our setting, if the governance condition does not reject updates in a neighborhood of \(x^\ast\), then the realized sequence \(x_t\) coincides with \(T\)-iteration near the fixed point and thus converges to the same \(x^\ast\).
\end{proof}

\begin{corollary}[Drift Vanishing at the Fixed Point]
\label{cor:drift-vanish}
Under the conditions of Theorem~\ref{thm:fixed-point}, the drift \(\delta_{\mathrm{drift}}(x_t, x_{t+1})\) converges to zero as \(t \to \infty\).
\end{corollary}

\begin{proof}
Since \(x_t \to x^\ast\) and \(T\) is continuous, \(x_{t+1} = T(x_t) \to T(x^\ast) = x^\ast\).%
Thus \(\|x_{t+1} - x_t\| \to 0\).%
By Lemma~\ref{lem:drift-global-norm},
\[
  \delta_{\mathrm{drift}}(x_t, x_{t+1})
  \le
  \|D\|\;\|x_{t+1} - x_t\| \to 0.
\]
\end{proof}

\subsection{Prime-Indexed Orchestration Operator Norms}

We now consider the prime-indexed orchestration structure.%
For simplicity, assume each component space \(S_p\) is finite-dimensional with norm \(\|\cdot\|_p\) induced by an inner product, so operator norms are well defined and equivalent across reasonable choices.

\begin{definition}[Local Update Operators]
For each prime \(p \in P\), let \(U_p : S_p \to S_p\) be a bounded linear operator with norm
\[
  \|U_p\| \coloneqq \sup_{\|z\|_p \le 1} \|U_p(z)\|_p.
\]
Define the block-diagonal operator \(U_{\Pi} : \prod_{p \in P} S_p \to \prod_{p \in P} S_p\) by
\[
  U_{\Pi}\bigl((x_p)_{p \in P}\bigr) = \bigl(U_p(x_p)\bigr)_{p \in P}.
\]
\end{definition}

\begin{lemma}[Norm of Block-Diagonal Prime Operator]
\label{lem:block-diagonal}
If \(\|U_p\| \le L_p\) for all \(p \in P\) and we set \(L \coloneqq \sup_{p \in P} L_p\), then
\[
  \|U_{\Pi}\| \le L
\]
with operator norm taken relative to \(\|\cdot\|_{\Pi}\).
\end{lemma}

\begin{proof}
Let \(x = (x_p)_{p \in P}\). Then
\begin{align*}
  \|U_{\Pi}(x)\|_{\Pi}^2
  &= \sum_{p \in P} \|U_p(x_p)\|_p^2
  \le \sum_{p \in P} L_p^2 \|x_p\|_p^2
  \le L^2 \sum_{p \in P} \|x_p\|_p^2
  = L^2 \|x\|_{\Pi}^2.
\end{align*}
Taking square roots yields \(\|U_{\Pi}(x)\|_{\Pi} \le L \|x\|_{\Pi}\).%
Thus \(\|U_{\Pi}\| \le L\).
\end{proof}

We can now bound the norm of the global update \(U : S \to S\) built from local updates via
\[
  U = R \circ U_{\Pi} \circ D.
\]

\begin{proposition}[Global Operator Norm via Prime Components]
\label{prop:global-norm}
Let \(U = R \circ U_{\Pi} \circ D\) as above, with \(\|D\|\), \(\|R\|\), and \(\|U_{\Pi}\|\) finite. Then
\[
  \|U\| \le \|R\|\,\|U_{\Pi}\|\,\|D\|.
\]
If \(\|U_{\Pi}\| \le L < \infty\), then \(\|U\| \le L\,\|D\|\,\|R\|\).
\end{proposition}

\begin{proof}
For any \(x \in S\),
\[
  \|U(x)\|
  =
  \|R(U_{\Pi}(D(x)))\|
  \le
  \|R\|\;\|U_{\Pi}(D(x))\|_{\Pi}
  \le
  \|R\|\;\|U_{\Pi}\|\;\|D(x)\|_{\Pi}
  \le
  \|R\|\;\|U_{\Pi}\|\;\|D\|\;\|x\|.
\]
Taking the supremum over \(\|x\| \le 1\) yields the result.
\end{proof}

\begin{corollary}[Contractiveness Condition via Prime Components]
\label{cor:contractive-prime}
If \(\|R\|\,\|U_{\Pi}\|\,\|D\| < 1\), then the global operator \(U\) is a contraction on \(S\).
\end{corollary}

\begin{proof}
Immediate from Proposition~\ref{prop:global-norm} and the definition of contraction.
\end{proof}

This provides a structural condition under which the prime-indexed local dynamics guarantee a contractive global update, thereby enabling fixed-point existence and convergence via Theorem~\ref{thm:fixed-point}.

\subsection{Governance Monotonicity and Lyapunov-type Bounds}

We now provide a more explicit Lyapunov-style argument for the governance functional \(G\).

\begin{assumption}[Governance Lyapunov Condition]
\label{assumption:governance-lyapunov}
There exists a function \(G : S \to \mathbb{R}\) such that:
\begin{enumerate}
  \item \(G(x) \ge 0\) for all \(x \in S\).
  \item \(G(x) = 0\) if and only if \(x \in \mathcal{A}\) (or a distinguished subset of \(\mathcal{A}\)).
  \item There exists \(\beta > 0\) such that for all accepted updates,
  \[
    G(x_{t+1}) \le G(x_t) - \beta\,\|x_{t+1} - x_t\|^2.
  \]
\end{enumerate}
\end{assumption}

\begin{theorem}[Governance as a Lyapunov Functional]
\label{thm:lyapunov}
Under Assumption~\ref{assumption:governance-lyapunov}, for any sequence \(\{x_t\}\) generated by the accept/reject rule:
\begin{enumerate}
  \item \(G(x_t)\) is nonincreasing and bounded below by \(0\), hence convergent.
  \item The series \(\sum_{t=0}^\infty \|x_{t+1} - x_t\|^2\) converges.
  \item \(\|x_{t+1} - x_t\| \to 0\), and any limit point of \(\{x_t\}\) lies in the set where \(G(x) = 0\), i.e., in \(\mathcal{A}\) (or the designated subset).
\end{enumerate}
\end{theorem}

\begin{proof}
By the update rule and Assumption~\ref{assumption:governance-lyapunov},
\[
  G(x_{t+1}) \le G(x_t) \quad \text{for accepted updates},
\]
and \(G(x_{t+1}) = G(x_t)\) when the update is rejected and \(x_{t+1} = x_t\).%
Thus \(G(x_t)\) is nonincreasing.%
Since \(G(x_t)\ge 0\), it converges to some \(G_\infty \ge 0\).

Furthermore, for each accepted update,
\[
  G(x_t) - G(x_{t+1}) \ge \beta \|x_{t+1} - x_t\|^2.
\]
Summing over \(t\) yields
\[
  \sum_{t=0}^\infty \|x_{t+1} - x_t\|^2
  \le
  \frac{1}{\beta} \bigl(G(x_0) - \lim_{T\to\infty} G(x_T)\bigr)
  \le
  \frac{G(x_0)}{\beta} < \infty.
\]
Therefore the series converges, and in particular \(\|x_{t+1} - x_t\| \to 0\).%

Finally, any limit point \(x_\infty\) of \(\{x_t\}\) must satisfy \(G(x_\infty) = G_\infty\).%
If \(G\) distinguishes lawful states by \(G(x) = 0\) iff \(x \in \mathcal{A}\), then \(G_\infty = 0\) implies \(x_\infty \in \mathcal{A}\).%
Even if \(G_\infty > 0\), any stable limit point must satisfy the condition that no further lawfulness-preserving, governance-decreasing update is possible, which in many designs corresponds to an attractor inside \(\mathcal{A}\).%
\end{proof}

Combining Theorem~\ref{thm:lyapunov} with Lemma~\ref{lem:drift-global-norm} yields a corresponding statement for the drift in the prime-indexed space:
\[
  \delta_{\mathrm{drift}}(x_t, x_{t+1})
  \le \|D\|\,\|x_{t+1} - x_t\| \to 0.
\]

\subsection{Summary of Operator-Norm Conditions}

For convenience, we collect the main operator-norm conditions that guarantee:
\begin{itemize}
  \item bounded drift,
  \item stability of recursion,
  \item convergence to a lawful attractor.
\end{itemize}

\begin{enumerate}
  \item \(\|D\| < \infty\), \(\|R\| < \infty\).
  \item The prime-indexed local operators \(U_p\) are bounded with \(\sup_p \|U_p\| < \infty\), so that \(\|U_{\Pi}\| < \infty\).
  \item The composite global operator \(U = R \circ U_{\Pi} \circ D\) satisfies
  \[
    \|U\| \le \|R\|\,\|U_{\Pi}\|\,\|D\| < 1
  \]
  (or, more generally, the \(\Pi_{\mathcal{A}} \circ U\) composite is contractive on \(\mathcal{A}\)).
  \item The correction operator \(\Pi_{\mathcal{A}}\) is nonexpansive (metric projection or equivalent), with \(\Pi_{\mathcal{A}}(x) = x\) for \(x \in \mathcal{A}\).
  \item The governance functional \(G\) satisfies the Lyapunov-like condition in Assumption~\ref{assumption:governance-lyapunov}.
\end{enumerate}

Under these conditions, the GRIMS+/prime-indexed recursive system admits:
\begin{itemize}
  \item A unique lawful fixed point (or attractor set) in \(\mathcal{A}\).
  \item Convergence to that point under the recursive accept/reject dynamics.
  \item Vanishing drift in the prime-indexed decomposition, ensuring prime-lawfulness equivalence in the limit.
\end{itemize}

These results provide a mathematically explicit backbone for the informal theorems articulated in the main article, and they show precisely how the prime-indexed, ethics-constrained recursion can be made stable and well-behaved in operator-norm terms.

\nocite{*}
\bibliographystyle{unsrt}  
\bibliography{references}
\end{document}